\section{The Original Pipeline and Its Audit}
\label{sec:original}

\subsection{Architecture}
The original code (preserved verbatim in \code{legacy/}) has three modules. Figure~\ref{fig:pipeline-old}
shows the data flow.

\begin{figure}[H]
\centering
\begin{tikzpicture}[node distance=0.35cm and 0.45cm,
  box/.style={draw,rounded corners,fill=white,align=center,font=\scriptsize,minimum height=0.95cm,text width=2.05cm},
  bad/.style={box,fill=red!8,draw=red!60!black},
  arr/.style={-{Stealth},thick}]
\node[box] (csv) {\code{output.csv}\\75k rows};
\node[bad,right=of csv] (clean) {clean labels,\\inf$\to$0\\\textit{keep duplicates}};
\node[bad,right=of clean] (mm) {MinMax$\to[0,\pi]$\\on \textbf{all} data\\pad to 58};
\node[box,right=of mm] (bal) {balance\\(unseeded)};
\node[box,right=of bal] (split) {80/20\\stratified split};
\node[box,below=0.9cm of split] (pca) {PCA$\to n$\\(train)\\min--max $[0,\pi]$};
\node[box,left=of pca] (zz) {ZZ map\\$\alpha=1$, $r$ reps\\Cirq};
\node[bad,left=of zz] (ker) {$K_{ij}$: 2 circuit\\sims per entry\\$O(N^2)$};
\node[box,left=of ker] (svc) {SVC(precomp.)\\$C=1$};
\node[bad,left=of svc] (rbf) {RBF-SVM\\$C{=}1,\gamma$=scale\\(untuned)};
\draw[arr] (csv)--(clean); \draw[arr] (clean)--(mm); \draw[arr] (mm)--(bal); \draw[arr] (bal)--(split);
\draw[arr] (split)--(pca); \draw[arr] (pca)--(zz); \draw[arr] (zz)--(ker); \draw[arr] (ker)--(svc);
\draw[arr,dashed] (pca.south) -- ++(0,-0.35) -| node[pos=0.25,below,font=\scriptsize]{same PCA features} (rbf.south);
\end{tikzpicture}
\caption{Original pipeline. Red boxes mark steps changed by this work.}
\label{fig:pipeline-old}
\end{figure}

\paragraph{\code{data\_loader.py}.} \code{load\_csv()} reads the CSV, strips column names, removes
repeated header rows, and draws up to $3\times$\code{max\_samples} rows at random
(\code{random\_state=42}). It builds binary or multiclass labels, selects up to 61 named feature
columns, replaces $\pm\infty$/NaN by 0, drops zero-variance columns, and pads or trims to exactly 58
columns. It then undersamples each class to $\min(\text{smallest class},\,\text{max\_samples}/k)$
using the \emph{unseeded} global \code{np.random}. Finally it fits a \code{MinMaxScaler} to $[0,\pi]$
on \emph{all} remaining rows and returns them.

\paragraph{\code{qsvm.py}.} \code{ZZFeatureMap.build\_encoding\_circuit()} builds the circuit of
Figure~\ref{fig:circuit} in Cirq. \code{build\_kernel\_circuit()} builds $U^\dagger(\bm a)U(\bm b)$
plus measurement, for the sampling mode. \code{QuantumKernel.\_kernel\_statevector(a,b)} simulates
\emph{both} circuits with \code{cirq.Simulator} and returns $|\langle\psi_a|\psi_b\rangle|^2$.
\code{compute\_kernel\_matrix()} loops over all $(i,j)$ pairs, using symmetry for the training
matrix. \code{QSVM.fit()} fits PCA to \code{n\_qubits} components on the training rows, rescales
each component to $[0,\pi]$ with training min/max, computes $K_{\text{train}}$ and fits
\code{SVC(kernel='precomputed', C=1, probability=True)}. \code{predict()} and
\code{predict\_proba()} \emph{each} rebuild the full test kernel.

\paragraph{\code{main.py}.} \code{run\_pipeline()} loads the data, splits it 80/20 (stratified,
\code{random\_state=42}), truncates the training set to 3000 rows, trains the QSVM, then trains
\code{SVC(kernel='rbf', C=1, gamma='scale')} on standardised PCA features (same projection) and
prints accuracy, a classification report, timings and a 4-panel figure (two confusion matrices,
ROC, accuracy bars).

\subsection{Reported results}
Table~\ref{tab:paper} reproduces the paper's headline numbers. Only the ZZ method exists in the
code. The Analytical and Sampling variants were run from earlier versions that are not in the
repository.

\begin{table}[H]
\centering\small
\caption{Results reported in the paper~\cite{ourpaper} (single run each).}
\label{tab:paper}
\begin{tabular}{llccc c}
\toprule
Method & Mode & QSVM acc. & SVM acc. & QSVM AUC & $N_\text{test}$ \\
\midrule
Analytical ($R_x$, 16q) & binary & 97.57\% & 98.54\% & 0.987 & 1028 \\
Analytical ($R_x$, 16q) & multiclass & 89.63\% & 87.68\% & -- & 3084 \\
Sampling (1024 shots)   & binary & 98.54\% & 98.54\% & 0.992 & 1028 \\
Sampling (1024 shots)   & multiclass & 89.63\% & 87.68\% & -- & 3084 \\
ZZ (8q, 2 reps)         & binary & \textbf{98.79\%} & 96.97\% & 0.999 & 822 \\
ZZ (16q, 1 rep)         & multiclass & \textbf{92.00\%} & 80.00\% & -- & 996 \\
\bottomrule
\end{tabular}
\end{table}

\subsection{Audit findings}
We read every line of the three modules and checked the data file. Table~\ref{tab:audit} lists
the issues found, ordered roughly by how much they can change the conclusions.

\begin{table}[H]
\centering\small
\caption{Issues found in the original pipeline.}
\label{tab:audit}
\begin{tabularx}{\linewidth}{c l X c}
\toprule
\# & Issue & Effect & Severity \\
\midrule
I1 & Exact duplicates kept & 65.9\% of rows are duplicates; $\approx$30\% of test flows have a
  training twin (measured). Accuracy is inflated by memorisation. & High \\
I2 & Untuned, unequal baselines & QSVM and RBF both fixed at $C=1$; RBF $\gamma$ fixed at
  \code{'scale'}. The comparison measures default settings, not models. & High \\
I3 & Kernel concentration & 16-qubit ZZ kernel with $\alpha=1$: median off-diagonal entry
  $\approx 2\times10^{-5}$, 84--89\% of entries below $10^{-3}$, so $K\approx I$ (measured in Section~\ref{sec:results}). No bandwidth control. & High \\
I4 & Cost $2N^2$ simulations & Two full circuit simulations per kernel entry, and the test kernel
  is computed twice. A 1600-sample, 16-qubit run is estimated at $\approx$44\,h. & High (practical) \\
I5 & Global MinMax before split & Test-set min/max leak into the training features. & Medium \\
I6 & Linear encoding of heavy tails & For most features, 90\% of flows occupy $<5\%$ of
  $[0,\pi]$, so qubits receive near-constant angles. & Medium \\
I7 & Single run, no significance & Differences of 1--2 points on 400--1000 test flows are within
  sampling noise. There are no error bars and no tests. & Medium \\
I8 & Non-reproducible sampling & Class balancing uses unseeded \code{np.random}; Platt scaling
  is unseeded. Each run draws different data. & Medium \\
I9 & Inconsistent sample sizes & The paper compares methods on different $N$ (822 vs.\ 1028 test). & Medium \\
I10 & Missing methods & The Analytical and Sampling kernels in the paper are not in the code. & Medium \\
I11 & Narrow metrics & Accuracy only. IDS operators care about false-alarm rate and detection
  rate. Multiclass AUC is not reported. & Low \\
I12 & Cosmetic & Zero-padding to 58 columns; hard-coded ``58 $\to$'' message; silent truncation
  to 3000 training rows. & Low \\
I13 & Pre-sampling shrinks the data & $3\times$\code{max\_samples} random rows are drawn \emph{before}
  balancing. BENIGN is only 6.8\% of the file, so \code{--samples 2000} yields 411 BENIGN rows and
  822 binary samples in total (165 test flows). & Medium \\
\bottomrule
\end{tabularx}
\end{table}

\paragraph{How small the binary test set was.} Running the original loader confirms I13: with the
default \code{--samples 2000} it returns exactly 822 binary samples, the $N$ quoted in the paper for
the ZZ binary experiment, so the test split had 165 flows. The reported 98.79\% vs.\ 96.97\% is then
$163/165$ vs.\ $160/165$, a difference of \textbf{three test flows}. Even if all three were cases
where only the QSVM was right, the exact McNemar test would give $p=0.25$. The multiclass figures
(92.00\% vs.\ 80.00\%, $N=996=6\times166$) correspond to a 200-flow test set. The 12-point gap is
large, but it comes from a single split with the issues above.

None of these issues is unusual. Duplicate leakage in particular is a well-documented trap in
intrusion-detection benchmarks~\cite{engelen2021,kaufman2012}, and untuned baselines are the most
common cause of overstated quantum advantage~\cite{bowles2024}. The core idea and the quantum
code itself are sound. The fast engine reproduces the original kernel values exactly
(Section~\ref{sec:verify}).
